\section{Related Work}
\label{sec:related-work}

\subsection{Knowledge tracing and student performance prediction}

Knowledge tracing (KT) estimates a student's evolving knowledge from a sequence
of learning interactions.  Bayesian Knowledge Tracing represents mastery of a
knowledge component with a hidden Markov model \citep{corbett1995knowledge},
while Performance Factors Analysis uses counts of prior successes and failures
as interpretable predictors \citep{pavlik2009performance}.  Item-response and
factor models instead explain outcomes through student ability and item
difficulty \citep{rasch1960probabilistic,vie2019knowledge}.  Neural KT replaces
hand-specified transitions with recurrent networks \citep{piech2015deep},
external memories \citep{zhang2017dynamic}, or self-attention
\citep{pandey2019self,ghosh2020context}.  These models are strong and appropriate
baselines for correctness prediction, but their standard objective and output
space remain task-specific: they estimate a binary outcome for an indexed item
or knowledge component.  Our work does not argue that such models are incapable
of richer prediction after modification.  Instead, we ask whether correctness,
concrete responses, misconception text, and prospective responses can be
handled by one generative interface without adding a separate prediction head
for every target.

The distinction also motivates our evaluation design.  Standard KT benchmarks
often emphasize warm-item, new-student prediction, where test students are held
out but most test items occurred during training.  Text-rich datasets such as
XES3G5M \citep{liu2023xes3g5m}, DBE-KT22
\citep{abdelrahman2022dbekt22}, and FoundationalASSIST
\citep{worden2026foundationalassist} permit models to use problem content in
addition to identifiers.  We retain student-disjoint temporal tests but report
protocols by data source and prediction type rather than pooling unlike label
spaces into a macro-average.

\subsection{Responses, misconceptions, and cognitive diagnosis}

Correctness does not identify the process that generated an answer.  Cognitive
diagnosis models infer mastery profiles over a predefined skill matrix
\citep{delatorre2009dina}, while distractor and misconception modeling links
particular wrong answers to systematic gaps in understanding.  The Diagnostic
Questions/Eedi benchmark made this connection explicit by associating
mathematics distractors with expert-authored misconceptions
\citep{wang2021diagnostic}.  More recently, MalruleLib formalized 101 executable
mathematical error rules and their step traces \citep{chen2026malrulelib}.  These
resources support complementary forms of evidence: Eedi provides authentic
expert labels for observed distractors, whereas MalruleLib provides controlled
counterfactual behavior and exact rule identity.

Our misconception task differs from ordinary per-example candidate ranking.
The main protocol either retrieves from the complete misconception taxonomy or
generates a diagnosis without exposing a candidate set constructed to contain
the gold label.  Our response tasks similarly distinguish closed-choice
prediction from open mathematical response generation.  This separation is
important: retrieving a supplied misconception and generating a student's
future response are not interchangeable capabilities, even when both can be
scored with semantic similarity.

\subsection{Language models as knowledge tracers and simulated students}

Recent work incorporates pretrained language models into KT to represent
question semantics and address cold-start items.  Language Model Can Do
Knowledge Tracing adapts language-model representations to KT
\citep{lee2024language}; CLST aligns a generative language model with a student
tracer for cold-start mitigation \citep{jung2024clst}; and SINKT combines
language-model semantics with structural information for inductive KT
\citep{fu2024sinkt}.  This line of work primarily retains correctness as the
prediction target.  It demonstrates the value of language semantics for KT but
does not by itself unify correctness with open responses and misconception
diagnosis.

A related line asks LLMs to act as students.  Sonkar and Baraniuk
\citeyearpar{sonkar2023deduction} probe whether LLMs follow provided student
evidence under perturbations.  Agent-based approaches condition LLMs on coarse
cognitive levels to generate imperfect behavior \citep{wu2025embracing}.
Contemporaneous work emphasizes that plausible generations are insufficient
evidence of valid student simulation and calls for comparisons against real
student behavior and stronger controls \citep{yuan2026valid}.  We share this
concern but study a different setting: the model is trained over actual student
histories and multiple predictive targets, rather than prompted to enact a
persona or ability level.  We also test whether history-dependent information
is localized in an internal state that can be intervened on.

\subsection{Probing and intervening on internal representations}

Linear probes are widely used to identify information encoded in neural
representations, but probe accuracy can reflect information that the model does
not use for its prediction.  Control tasks and selectivity analyses help
separate representation content from probe capacity \citep{hewitt2019designing},
while counterfactual removal or intervention tests whether an encoded feature is
behaviorally relevant \citep{elazar2021amnesic,ravfogel2021counterfactual}.
Accordingly, we do not treat a decodable student attribute as sufficient
evidence for a functional student state.

Our architecture goes beyond extracting a token activation from a complete
history.  It recurrently updates a compact continuous state and removes the raw
history from the prediction path.  Full-context distillation tests whether this
state preserves behaviorally useful information, while future-profile
supervision ties the state to later student outcomes.  At evaluation, a frozen
readout receives no history or student ID; permuting state ownership across
held-out students therefore tests whether the student-specific representation
matters.  Disjoint-history views test within-student stability, and population
and same-window outcome baselines test incremental utility.  We separately
report donor, same-profile, random, and no-op text-transport interventions as
diagnostics, without requiring arbitrary cross-student latent transplants to be
in distribution.
